% Source: OpenAI2026AreaLaw, 09-amplification.tex, lines 101--139 and 173,
% openai/math@adc7f1241b42e322a6451854ab7e4b4c146bf78a.
% Checked declarations: LocalChannelOscillation at 436daa61 and
% ExponentialChannelShells at 73381fcb; SpectatorOscillationSupport at 54772a7c.
% GraphChannelOscillation production passed native compilation at 02ec325d.
% This leaf covers initial support, finite-radius physical shells, and
% one-event increments, including full channels with exact component support.
% It does not assert an infinite-radius limit or propagation under
% retained Poisson clocks.

\section{Finite one-event oscillation bounds}
\label{sec:al_channel_propagation}

Use the physical space $\mathcal H=(\C^q)^{\otimes I}$, with $q>0$,
and the arbitrary finite spectator $\mathcal K=\C^J$ from
Section~\ref{sec:al_channel_shell_oscillation}; the dimension of
$\mathcal K$ may be zero. Fix a positive contraction $0\le k\le\mathbb 1$
on $\mathcal H$ and increasing regions $(K_l)_{l\ge0}$, so
$K_l\subseteq K_m$ whenever $l\le m$. Write
$\mathcal E_l=\mathcal E_{K_l}$ for the localized root channel of
Definition~\ref{def:al_local_root_channel}, and $\Delta_l$ for its
shells from Definition~\ref{def:al_local_channel_shell}.
For every site $y\in I$, use the oscillation $d_y$ of
Definition~\ref{def:qc_site_oscillation}.
Throughout, $B$ is an arbitrary complex operator on
$\mathcal H\otimes\mathcal K$, without a Hermiticity assumption.
All norms are Euclidean operator norms.

\subsection{Initial oscillation from physical support}

For a physical region $S\subseteq I$, write $\mathcal A_S$ for the
physical operators supported on $S$. Let $E_S^{\mathcal K}$ be the
expectation of Definition~\ref{def:qc_spectator_expectation}, which
retains $S$ and leaves both spectator indices unchanged.
For $a,b\in J$, write $B_{ab}$ for the physical matrix with entries
$(B_{ab})_{\sigma,\tau}=B_{(\sigma,a),(\tau,b)}$.

\begin{theorem}[Physical support and spectator expectation]
    \label{thm:al_spectator_support_fixed}
    \uses{def:qc_spectator_expectation,def:qc_local_circuit}
    \lean{QuantumCircuit.spectatorSiteExpectation_eq_self_iff}
    \leanok
    For every region $S\subseteq I$,
    \begin{align}
        E_S^{\mathcal K}(B)=B
         & \quad\Longleftrightarrow\quad
        \forall a,b\in J,\ B_{ab}\in\mathcal A_S.
        \label{eq:al_propagation_support_fixed}
    \end{align}
    This equivalence remains valid for matrices with an arbitrary
    spectator index set $J$.
\end{theorem}
\begin{proof}
    \leanok
    \uses{def:qc_spectator_expectation,lem:qc_site_expectation_properties}
    By~(\ref{eq:qc_spectator_blocks}), the fixed-point condition is
    equivalent to $E_S(B_{ab})=B_{ab}$ for every $a,b\in J$.
    Each $E_S(B_{ab})$ belongs to $\mathcal A_S$, proving the forward
    implication. Conversely, $E_S$ fixes every matrix in
    $\mathcal A_S$, so equality holds in every matrix entry.
\end{proof}

\begin{theorem}[Vanishing oscillation outside the physical support]
    \label{thm:al_spectator_support_oscillation_zero}
    \uses{def:qc_site_oscillation,def:qc_local_circuit}
    \lean{QuantumCircuit.siteOscillation_eq_zero_of_not_mem_support}
    \leanok
    Suppose $B_{ab}\in\mathcal A_S$ for every $a,b\in J$.
    If $y\notin S$, then $d_y(B)=0$.
\end{theorem}
\begin{proof}
    \leanok
    \uses{thm:al_spectator_support_fixed,thm:qc_spectator_commutant,
        thm:qc_site_oscillation_commutant}
    Theorem~\ref{thm:al_spectator_support_fixed} gives
    $E_S^{\mathcal K}(B)=B$. Every $U\in\mathcal A_{\{y\}}$
    belongs to $\mathcal A_{I\setminus S}$, hence
    $[B,\widehat U]=[E_S^{\mathcal K}(B),\widehat U]=0$ by
    Theorem~\ref{thm:qc_spectator_commutant}.
    Theorem~\ref{thm:qc_site_oscillation_commutant} now gives
    $d_y(B)=0$.
\end{proof}

\begin{theorem}[Initial oscillation bounded by the support indicator]
    \label{thm:al_spectator_support_indicator}
    \uses{def:qc_site_oscillation,def:qc_local_circuit}
    \lean{QuantumCircuit.siteOscillation_le_support_indicator}
    \leanok
    Suppose $B_{ab}\in\mathcal A_S$ for every $a,b\in J$.
    For every $y\in I$,
    \begin{align}
        d_y(B) & \le2\|B\|\mathbf1_S(y),
        \label{eq:al_propagation_support_indicator}
    \end{align}
    where $\mathbf1_S(y)$ is one on $S$ and zero outside $S$.
    This is the initial support estimate in
    \cite[Section~10]{OpenAI2026AreaLaw}, for arbitrary finite
    spectators and arbitrary complex $B$.
\end{theorem}
\begin{proof}
    \leanok
    \uses{thm:al_spectator_support_oscillation_zero,
        thm:qc_site_oscillation_bounds}
    If $y\in S$, apply $d_y(B)\le2\|B\|$ from
    Theorem~\ref{thm:qc_site_oscillation_bounds}.
    If $y\notin S$, both sides of
    (\ref{eq:al_propagation_support_indicator}) are zero by
    Theorem~\ref{thm:al_spectator_support_oscillation_zero}.
\end{proof}

\subsection{Coefficients from the localization errors}

\begin{definition}[Shell coefficients from localization errors]
    \label{def:al_channel_shell_coefficient}
    \lean{TNLean.PEPS.AreaLaw.localRootChannelShellBound}
    \leanok
    For a real sequence $(\varepsilon_l)_{l\ge0}$, define
    \begin{align}
        \eta_0 & =2,
        \label{eq:al_propagation_coefficient_zero}
    \end{align}
    \begin{align}
        \eta_{l+1} & =4\bigl(\sqrt{\varepsilon_{l+1}}
        +\sqrt{\varepsilon_l}\bigr).
        \label{eq:al_propagation_coefficient_successor}
    \end{align}
    Here the square root of a negative real number is interpreted as
    zero. Every localization-error inequality below forces its error
    bound to be nonnegative at the radius in question.
\end{definition}

\begin{theorem}[Nonnegativity of the shell coefficients]
    \label{thm:al_channel_shell_coefficient_nonneg}
    \uses{def:al_channel_shell_coefficient}
    \lean{TNLean.PEPS.AreaLaw.localRootChannelShellBound_nonneg}
    \leanok
    For every real sequence $(\varepsilon_l)$ and every $l\ge0$,
    the coefficient $\eta_l$ is nonnegative.
\end{theorem}
\begin{proof}
    \leanok
    \uses{def:al_channel_shell_coefficient}
    The zeroth coefficient is two, and each successor coefficient
    is four times a sum of nonnegative square roots.
\end{proof}

\begin{theorem}[Finite-radius physical shell bound]
    \label{thm:al_channel_shell_finite_bound}
    \uses{def:al_local_channel_shell,def:qc_site_oscillation,
        def:al_channel_shell_coefficient}
    \lean{TNLean.PEPS.AreaLaw.norm_localRootChannelShell_le_sum_siteOscillation}
    \leanok
    Fix $l\in\N$ and a real sequence $(\varepsilon_j)$.
    Suppose that, for every integer $j$ with $0\le j\le l$,
    \begin{align}
        \|k-E_{K_j}(k)\| & \le\varepsilon_j.
        \label{eq:al_propagation_finite_error}
    \end{align}
    Then
    \begin{align}
        \|\Delta_l(B)\| & \le\eta_l\sum_{z\in K_l}d_z(B).
        \label{eq:al_propagation_finite_shell}
    \end{align}
    In particular, if $\|k-E_{K_j}(k)\|=0$ for every
    $0\le j\le L$, then $\Delta_j(B)=0$ for $1\le j\le L$.
    The zeroth shell still has coefficient $\eta_0=2$.
\end{theorem}
\begin{proof}
    \leanok
    \uses{thm:al_channel_shell_oscillation,def:al_channel_shell_coefficient}
    At $l=0$, apply~(\ref{eq:al_shell_zero_oscillation}).
    At $l=m+1$, use the error inequalities at $m+1$ and $m$ in
    (\ref{eq:al_shell_successor_oscillation}), with coefficient
    (\ref{eq:al_propagation_coefficient_successor}).
    When both errors are zero, this coefficient is zero, and the
    shell has norm zero. The definition of $\eta_0$ is independent
    of the errors.
\end{proof}

\begin{theorem}[Finite one-event oscillation increment]
    \label{thm:al_channel_finite_increment}
    \uses{def:al_local_root_channel,def:qc_site_oscillation,
        def:al_channel_shell_coefficient}
    \lean{TNLean.PEPS.AreaLaw.siteOscillation_localRootChannel_le}
    \leanok
    Fix $n\in\N$, $y\in I$, and a real sequence $(\varepsilon_l)$.
    Suppose that $\|k-E_{K_l}(k)\|\le\varepsilon_l$ for every
    integer $l$ with $0\le l\le n$. Then
    \begin{align}
        \label{eq:al_propagation_finite_increment}
        d_y(\mathcal E_n(B))
         & \le d_y(B)+2\sum_{\substack{0\le l\le n \\y\in K_l}}
        \eta_l\sum_{z\in K_l}d_z(B).
    \end{align}
    Every regional sum is evaluated on the original operator $B$.
    This is the finite-radius one-event estimate underlying
    \cite[Section~10]{OpenAI2026AreaLaw}.
\end{theorem}
\begin{proof}
    \leanok
    \uses{thm:al_channel_shell_finite_bound,
        thm:al_channel_shell_coefficient_nonneg,
        thm:al_root_channel_outside_oscillation,
        thm:qc_site_oscillation_seminorm,
        thm:qc_site_oscillation_bounds,def:al_local_channel_shell}
    Put $a_l=\eta_l\sum_{z\in K_l}d_z(B)$, which is nonnegative.
    Induct on $n$. If $y\notin K_n$, then
    $d_y(\mathcal E_n(B))\le d_y(B)$ by
    Theorem~\ref{thm:al_root_channel_outside_oscillation}, and the
    nonnegative sum in~(\ref{eq:al_propagation_finite_increment})
    preserves this inequality.
    At $n=0$ with $y\in K_0$, write
    $\mathcal E_0(B)=B+\Delta_0(B)$ to obtain
    \begin{align}
        d_y(\mathcal E_0(B))
         & \le d_y(B)+d_y(\Delta_0(B))\le d_y(B)+2a_0.
        \label{eq:al_propagation_increment_zero}
    \end{align}
    For the successor step with $y\in K_{n+1}$, the identity
    $\mathcal E_{n+1}(B)=\mathcal E_n(B)+\Delta_{n+1}(B)$ gives
    \begin{align}
        d_y(\mathcal E_{n+1}(B))
         & \le d_y(\mathcal E_n(B))+2\|\Delta_{n+1}(B)\|
        \le d_y(\mathcal E_n(B))+2a_{n+1}.
        \label{eq:al_propagation_increment_step}
    \end{align}
    Restrict the error assumptions to radii at most $n$ and apply
    the induction hypothesis to the first term in
    (\ref{eq:al_propagation_increment_step}). The new term is
    exactly the contribution of $n+1$ to the filtered sum.
\end{proof}

\subsection{Stretched-exponential localization errors}

\begin{theorem}[Explicit stretched-exponential shell coefficients]
    \label{thm:al_channel_shell_coefficient_exp}
    \uses{def:al_channel_shell_coefficient}
    \lean{TNLean.PEPS.AreaLaw.localRootChannelShellBound_le_exp}
    \leanok
    Let $C\ge0$, $c>0$, and $0<\alpha\le1$. Set
    $\varepsilon_l=C\exp(-cl^\alpha)$ and
    \begin{align}
        A & =2+8\sqrt C\exp(c/2), & b & =c/2.
        \label{eq:al_propagation_exp_constants}
    \end{align}
    For every $l\in\N$,
    \begin{align}
        \eta_l & \le A\exp(-bl^\alpha).
        \label{eq:al_propagation_exp_coefficient}
    \end{align}
\end{theorem}
\begin{proof}
    \leanok
    \uses{def:al_channel_shell_coefficient}
    The identity $\sqrt{C\exp(-cl^\alpha)}
        =\sqrt C\exp(-bl^\alpha)$ rewrites each successor coefficient.
    Since $0<\alpha\le1$, one has $(l+1)^\alpha\le l^\alpha+1$,
    and therefore
    \begin{align}
        \exp(-bl^\alpha)
         & \le\exp(b)\exp(-b(l+1)^\alpha).
        \label{eq:al_propagation_exp_previous}
    \end{align}
    Also $1\le\exp(b)$. Apply this inequality and
    (\ref{eq:al_propagation_exp_previous}) to the two square roots
    in~(\ref{eq:al_propagation_coefficient_successor}) to obtain
    \begin{align}
        \eta_{l+1}
         & \le8\sqrt C\exp(b)\exp(-b(l+1)^\alpha)
        \le A\exp(-b(l+1)^\alpha).
        \label{eq:al_propagation_exp_successor}
    \end{align}
    At radius zero, $0^\alpha=0$ and $\eta_0=2\le A$, proving
    (\ref{eq:al_propagation_exp_coefficient}) there as well.
\end{proof}

\begin{theorem}[Stretched-exponential physical shell bound]
    \label{thm:al_channel_shell_exp_bound}
    \uses{def:al_local_channel_shell,def:qc_site_oscillation}
    \lean{TNLean.PEPS.AreaLaw.norm_localRootChannelShell_le_exp_sum_siteOscillation}
    \leanok
    Let $C\ge0$, $c>0$, and $0<\alpha\le1$, and take $A,b$ from
    (\ref{eq:al_propagation_exp_constants}). Fix $l\in\N$.
    Suppose that, for every integer $j$ with $0\le j\le l$,
    \begin{align}
        \|k-E_{K_j}(k)\| & \le C\exp(-cj^\alpha).
        \label{eq:al_propagation_exp_error}
    \end{align}
    Then
    \begin{align}
        \|\Delta_l(B)\|
         & \le A\exp(-bl^\alpha)\sum_{z\in K_l}d_z(B).
        \label{eq:al_propagation_exp_shell}
    \end{align}
\end{theorem}
\begin{proof}
    \leanok
    \uses{thm:al_channel_shell_finite_bound,
        thm:al_channel_shell_coefficient_exp,thm:qc_site_oscillation_bounds}
    Apply~(\ref{eq:al_propagation_finite_shell}) with
    $\varepsilon_j=C\exp(-cj^\alpha)$. Bound its coefficient by
    (\ref{eq:al_propagation_exp_coefficient}) and use
    $\sum_{z\in K_l}d_z(B)\ge0$.
\end{proof}

\begin{theorem}[Stretched-exponential finite one-event increment]
    \label{thm:al_channel_exp_increment}
    \uses{def:al_local_root_channel,def:qc_site_oscillation}
    \lean{TNLean.PEPS.AreaLaw.siteOscillation_localRootChannel_le_exp}
    \leanok
    Let $C\ge0$, $c>0$, and $0<\alpha\le1$, and take $A,b$ from
    (\ref{eq:al_propagation_exp_constants}). Fix $n\in\N$ and
    $y\in I$. Suppose that
    $\|k-E_{K_l}(k)\|\le C\exp(-cl^\alpha)$ for every integer
    $l$ with $0\le l\le n$. Then
    \begin{align}
        \label{eq:al_propagation_exp_increment}
        d_y(\mathcal E_n(B))
         & \le d_y(B)+2\sum_{\substack{0\le l\le n \\y\in K_l}}
        A\exp(-bl^\alpha)\sum_{z\in K_l}d_z(B).
    \end{align}
    The estimate holds uniformly over all finite spectators and all
    complex operators $B$ under the stated finite set of error bounds.
\end{theorem}
\begin{proof}
    \leanok
    \uses{thm:al_channel_finite_increment,
        thm:al_channel_shell_coefficient_exp,thm:qc_site_oscillation_bounds}
    Apply~(\ref{eq:al_propagation_finite_increment}) with
    $\varepsilon_l=C\exp(-cl^\alpha)$. For each included radius,
    multiply~(\ref{eq:al_propagation_exp_coefficient}) by the
    nonnegative sum $\sum_{z\in K_l}d_z(B)$, and sum over exactly
    those radii with $0\le l\le n$ and $y\in K_l$.
\end{proof}

\subsection{The full channel on an anchor component}

Use the graph notation from Section~\ref{sec:qc_component_channel}:
$G$ is a simple graph on $I$, $a\in I$ is an anchor, $d_G$ is the
extended graph distance, $C(a)$ is the connected component of $a$,
and $N_l(a)=\{z\in I:d_G(a,z)\le l\}$.
Write $R_a$ for the maximum distance from $a$ to a vertex in $C(a)$.
Equivalently, maximize over all vertices after assigning zero distance
to unreachable vertices. Let $\mathcal E_k$ be the full root channel
in~(\ref{eq:qc_component_root_channel}).

\begin{theorem}[Full-channel increment with an explicit finite cutoff]
    \label{thm:al_graph_channel_exp_increment}
    \uses{def:al_local_root_channel,def:qc_site_oscillation,
        def:qc_local_circuit}
    \lean{TNLean.PEPS.AreaLaw.siteOscillation_spectatorRootChannel_le_exp_of_component_support}
    \leanok
    Suppose $0\le k\le\mathbb 1$ and $k\in\mathcal A_{C(a)}$.
    Let $C\ge0$, $c>0$, $0<\alpha\le1$, and take $A,b$ from
    (\ref{eq:al_propagation_exp_constants}). Fix an integer
    $n\ge R_a$. Suppose that
    $\|k-E_{N_l(a)}(k)\|\le C\exp(-cl^\alpha)$ for every integer
    $l$ with $0\le l\le n$. Then, for every $y\in I$,
    \begin{align}
        \label{eq:al_propagation_graph_increment}
        d_y(\mathcal E_k(B))
         & \le d_y(B)+2\sum_{\substack{0\le l\le n \\d_G(a,y)\le l}}
        A\exp(-bl^\alpha)\sum_{z\in N_l(a)}d_z(B).
    \end{align}
    All regional sums use the original operator $B$.
\end{theorem}
\begin{proof}
    \leanok
    \uses{thm:al_channel_exp_increment,thm:qc_graph_ball_stabilization,
        thm:qc_component_channel_stabilization}
    Apply Theorem~\ref{thm:al_channel_exp_increment} to the increasing
    regions $K_l=N_l(a)$. Since $n\ge R_a$ and
    $k\in\mathcal A_{C(a)}$, exact component localization gives
    $\mathcal E_{N_n(a)}(B)=\mathcal E_k(B)$ by
    Theorem~\ref{thm:qc_component_channel_stabilization}.
    Substitute this identity into~(\ref{eq:al_propagation_exp_increment})
    and rewrite $y\in N_l(a)$ as $d_G(a,y)\le l$.
\end{proof}
